\section{本章小结}
\label{sec:graph-sparsification-summary}

本章研究如何用边数少得多的带权图代替原图，同时近似保持计算结果。

若替代图只用于评价分组，只需保持割权——单点割给出加权度，加权度相加给出体积，割权与体积共同决定导通率和NCut。若还要求特征值、有效电阻或方程解，就需要保持所有实值顶点赋值上的能量，即谱近似。同一个三角形的例子说明，割近似不自动蕴含谱近似：边权的增减可能集中在差异大或差异小的方向上，只看跨边总权无法发现。

每条边在随机生成树中出现的概率是$c_eR_e$，这个乘积由电导和有效电阻共同决定。桥必定出现在每棵树中，团内边只出现在一部分树中。把这些概率除以Foster总和$n-1$，得到加起来等于一的抽样分布。每次抽到一条边后按概率的倒数补偿权重，使平均矩阵恰好等于原矩阵。有效电阻抽样的关键性质是：不论抽到哪条边，单次贡献在标准化坐标中的范数都是同一个值$(n-1)/q$，它均衡了所有边在最坏方向上的波动。矩阵Chernoff界把这个单次上界转为"所有方向同时满足误差"的概率保证。

用稀疏图作预条件时，每轮在原图上算残差，在稀疏图上根据残差求修正方向。两图能量之比越接近一，换坐标后的矩阵越接近恒等算子，条件数越小，所需迭代次数越少。实际节省多少工作，还要加上构造稀疏图和准备其求解器的开销。
